
\begin{frame}{Appendix --- Turnout Placebo (compulsory vote)}
\hypertarget{app:turnoutplacebo}{}
\small
Turnout is one act per voter, so it is office-invariant and reported here on its own. Because voting is compulsory in Brazil, overall turnout is institutionally floored and should not respond to litigation, which makes it a placebo. A near-zero, insignificant coefficient is what a valid design should deliver, since it shows the instrument is not picking up an omitted municipal shock that moves all electoral behavior at once. The elastic facultative turnout margins that can move are tested separately \hyperlink{src:turnoutprofile}{\beamergotobutton{by voter profile}}.
\smallskip
\begin{center}
\scriptsize
\resizebox{0.6\linewidth}{!}{% Auto-generated by code/03_estimation/02_iv_main.R
% Do not edit manually -- rerun the generating script to update

\begin{tabular}{lc}
   \tabularnewline \midrule \midrule
   Dependent Variable: & Turnout (any ballot)\\  
   Model:              & (1)\\  
   Judicialization     & -0.006\\   
                       & (0.004)\\   
   2024 Mean           & 0.835\\  
   2016 Mean           & 0.856\\  
   \midrule \midrule
   \multicolumn{2}{l}{\emph{Clustered (state) standard-errors in parentheses}}\\
   \multicolumn{2}{l}{\emph{Signif. Codes: ***: 0.01, **: 0.05, *: 0.1}}\\
\end{tabular}
}
\end{center}
\smallskip
\footnotesize The coefficient is indistinguishable from zero ($\TurnoutCoef$, $p=\TurnoutP$). Two caveats: turnout fails the pre-trend placebo ($p=\PtRFTurnoutP$), and compulsory turnout, measured separately, rises ($\CompulsoryCoef$, $p=\CompulsoryP$, tF interval excludes zero). Outcome is the 2024 turnout rate controlling for the 2016 level; state (UF) FE; SEs clustered by state.\hfill\hyperlink{src:voterballot}{\beamerreturnbutton{Back}}
\end{frame}
